% ===== §1 Phenomenology: the moment of the nameless =====
\section{The Phenomenology of Nameless Co-Generation}
\label{p10:sec:phenomenology}

\noindent
This paper takes as its primitive datum a structure of experience that precedes every theoretical apparatus and to which every such apparatus must answer: the joint apprehension, by two persons, of a value that is being generated between them and that is not yet exhausted by symbolisation. The value is, in the terms of \xref{p10:sec:registers}, an emergent of the coupled dynamics of the three registers, and the datum is its co-generation: a value with imaginary, symbolic, and real aspects, transversal across the three and resident in none. The structure is to be distinguished from three phenomena with which it is readily conflated, the agreement of two parties on a sign, the convergence of two gazes upon a shared object, and the exchange of a token whose content each party antecedently possesses. It is none of these, and the ground of the distinction is given in \xref{p10:sec:registers}: each of the three is the reduction of value to a single register, the suppression of a coupling, which extinguishes the emergent. What is apprehended is in the process of generation; it is not yet, and will never be, exhausted by its symbolisation; and it is apprehended by each party under the aspect of being co-apprehended, the reciprocal awareness entering constitutively into what is generated. The methodological order of the paper follows the order of grounding. The phenomenological description is given first and without explanatory supplement (\xref{p10:sec:phenomenology}, \xref{p10:sec:structure}); the ontology is then derived as the set of conditions under which the described structure is possible (\xref{p10:sec:ontology}); and the formal epistemologies (quantum-structural and political-economic) are introduced last, as the formal articulation of what the ontology establishes rather than as independent posits laid over the description (\xref{p10:sec:grammar}, \xref{p10:sec:political}). The motivation for this order is substantive rather than expository. Each formal apparatus is a settled symbolic system imported from a domain in which its results are already cashed; to introduce such an apparatus prior to the description is to expend, in advance, a determination the description has not yet underwritten, which is the structural form of what the prior paper of this series termed settling failure. The apparatus is accordingly earned from the phenomenon and not expended upon it.

\subsection{The datum described}
\label{p10:sec:bare}

The following case fixes the datum in its minimal form. It is offered as a description of a recurrent structure of experience, not as an anecdote.

\begin{casebox}{the unnamed value}
Two persons are engaged in unremarkable conversation, the arrangement of a shared meal, the route of an afternoon walk. At a determinate moment the register of the exchange alters without announcement. A silence supervenes which is not the cessation of speech but a positive datum: a salience in the interval between them, apprehended by each as not deposited there by either. One party may, at a later time, attempt a name for it, \emph{the beginning}; \emph{the point at which I knew}, but the name is supplied retrospectively, is acknowledged by both as inadequate, and is treated by both as a record of the event rather than the event. What is given in the moment itself is not the content of what has arrived but the fact of its arrival; that the arrival is simultaneous for both; and that each apprehends in the other's expression the same arrival, such that the co-apprehension is internal to what is apprehended.
\end{casebox}

\noindent
Reduced to its structure, the datum exhibits three features that the standard models of attention do not accommodate. The \emph{first} is that the term apprehended is not an object antecedently present. In the developmental account of joint attention, the infant and caregiver orienting jointly toward a ball or an indicated bird \parencite{trevarthen1979,zahavi1975}, there is a determinate third term, an object antecedently present in the world, upon which two lines of regard converge. The present structure has no such antecedent term. What is jointly apprehended is a value in the course of emergence on the coupled dynamics of the registers (\xref{p10:sec:registers}): a generation in progress, transversal across the imaginary, symbolic, and real, and localisable in none. It is not that the term cannot be symbolised, but that its symbolisation, where it occurs, does not exhaust it, the real aspect of the emergent exceeding any signifier that names it. The \emph{second} feature is that the apprehension is not a registration. Neither party detects a value antecedently present and awaiting detection; the value is not prior to the apprehension. The \emph{third} feature is that the reciprocity of awareness is constitutive rather than cumulative. The structure is not the conjunction of two registrations (one party's apprehension of the term and, separately, that party's apprehension that the other apprehends it). The reciprocal awareness is itself a component of what is generated and a condition of its generation.

\subsection{Joint attention at its limit, and the insufficiency of the cognitive-scientific model}
\label{p10:sec:limit}

The term \emph{joint attention} is retained, on two grounds. The lineage is pertinent: the prior treatment in this series established joint attention as the mechanism through which two consciousnesses are mutually constituted and the subject is first formed (\pinkref{Paper~V}). And the structure at issue is continuous with that mechanism, displaced to a new term. The displacement is nonetheless decisive. The joint attention of developmental psychology is triadic and objectual: two subjects, one shared object, the geometry of two regards meeting at a thing. The co-apprehension of value requires a joint attention whose vertex is not an object but an opening, a convergence not upon a thing but upon a site at which something is not yet determinate. The structure is therefore not that of developmental joint attention but is more adequately described by a phenomenology of intercorporeality, in which two embodied subjects are attuned prior to the availability of a concept \parencite{merleau2012}, and by the psychoanalytic account of an attention organised around a void rather than directed upon a thing.

The distinction is consequential and not terminological. The persistent error is the assimilation of the limit-structure to the standard case, the construal of the co-apprehension of value as the attainment, by two parties, of agreement upon a shared sign whose content is thereby settled. The error is not that the two parties symbolise the value; they do, and the making of a shared sign for it is, as \xref{p10:sec:registers} establishes and the account of language in this series develops, internal to the value's sustenance. The error is the construal of that symbolisation as a \emph{settling}: as the severance of the symbol from the Real, the cashing of the value as a determinate token whose content is exhausted in the signifier. Symbolisation of this kind suppresses the coupling of the Symbolic with the Real and extinguishes the emergent (\xref{p10:sec:emergence}); it is the operation the series identifies as settling failure. The contrary symbolisation, which sustains the coupling of the Symbolic with the Real, the poetic symbolisation of the later analysis (\xref{p10:sec:subtractive}), does not extinguish the value but is the mode in which it is jointly held and recurrently regenerated. The structure to be described is therefore not the absence of a sign but a symbolisation that sustains rather than severs the coupling to the Real, a shared sign whose content is not exhausted in its being named.

\subsection{The sense in which the value is nameless}
\label{p10:sec:nameless-corrected}

The locution that the value is nameless, employed throughout the description, is to be understood in the determinate sense the preceding establishes, and not in the sense its surface suggests. The value is not nameless in the sense of being unsymbolisable, an unspeakable remainder of the Real that any name would falsify; that construal reduces value to a single register and is excluded by \xref{p10:sec:emergence}. The value is nameless in the sense that its symbolisation, which is possible and is undertaken, does not exhaust it: the value being an emergent transversal across the three registers, its real aspect exceeds any signifier that names it, such that the name is at once adequate as a sign held jointly and inadequate as an exhaustive determination. To say that the value is nameless is to say that no name closes it, not that no name reaches it. The two persons name the value, make a sign for it, bring it into the Symbolic; and the value, named, remains in generation, because the naming that sustains the coupling to the Real does not terminate the emergent but is one of the operations by which it is sustained. The description that follows is to be read under this correction throughout.

\subsection{The constitutive identity of apprehension and generation}
\label{p10:sec:paradox}

The three features converge on a single structural result, which is to be stated as the result it is rather than dissolved. In the standard case of perception, apprehension and generation are distinct: the object is present independently of its being perceived, and the perceiving contributes nothing to the object. In the described structure they are not distinct. The co-apprehension of the value, in the sense of \xref{p10:sec:nameless-corrected}, is identically its generation. The two parties do not detect a value and subsequently attend to it; in the joint sustaining of an attention whose term is not exhausted by its symbolisation, the value is generated. Apprehension and generation are one act.

This result determines the conditions to be discharged by the ontology. A value generated in the act of being co-apprehended cannot be an entity antecedent to its apprehension; its mode of being is not that of an object. An apprehension that generates its own term cannot be a registration of antecedent states; its mode of cognition is not representational. The phenomenological description thus constrains both the ontology of value (value is non-antecedent) and the epistemology of its apprehension (cognition is non-representational), and it imposes these constraints prior to the introduction of any formal apparatus. The apparatus introduced in \xref{p10:sec:grammar} is accordingly the formal articulation of a structure the description has already established, and not an external warrant. The remainder of the paper is constrained by the datum fixed here: each subsequent construction is to be assessed by its adequacy to the co-apprehension of a value transversal across the registers, generated in the act of its joint recognition and not exhausted by its symbolisation, and to nothing of larger scope.
